\documentclass[11pt]{amsart}

\usepackage[T1]{fontenc}
\usepackage{lmodern}
\usepackage{microtype}
\usepackage{mathtools,amssymb,amsthm}
\usepackage{array}
\usepackage[hidelinks]{hyperref}

\hypersetup{
  pdftitle={A Proof of the Bhagat-Kulkarni-Larsson-Murali Smallest-Discrepancy-Heap Formula}
}

\setlength{\emergencystretch}{3em}

\newtheorem{theorem}{Theorem}
\newtheorem{lemma}[theorem]{Lemma}
\newtheorem{proposition}[theorem]{Proposition}
\newtheorem{corollary}[theorem]{Corollary}
\theoremstyle{definition}
\newtheorem{conjecture}{Conjecture}
\theoremstyle{remark}
\newtheorem*{remark}{Remark}

\DeclareMathOperator{\FvF}{FvF}
\DeclareMathOperator{\AvA}{AvA}

\title[The smallest discrepancy heap]
{A Proof of the Bhagat--Kulkarni--Larsson--Murali\\
 Smallest-Discrepancy-Heap Formula}

\date{}

\subjclass[2020]{91A46, 91A05, 91A10}
\keywords{cumulative subtraction game, self-interest game, pure subgame perfect
equilibrium, tie-breaking, friendly, antagonistic, discrepancy}

\begin{document}

\begin{abstract}
In the two-player self-interest cumulative subtraction game on the subtraction
set $S=\{s_2,s_1\}$, $0<s_2<s_1$, the mover takes some $s\in S$ with $s\le h$
from a heap of size $h$ and adds $s$ to their own score. Indifferent positions
must be resolved by a convention: under $\FvF$ the mover breaks a tie in the
opponent's favour, under $\AvA$ against it. Bhagat, Kulkarni, Larsson and
Murali conjectured that in the non-dominant regime $s_1<2s_2$ with
$s_1/s_2\neq(k+1)/k$ for every $k\in\mathbb N$ some heap distinguishes the two
conventions, and that the smallest such heap is
$h=(n+2)s_2+ns_1$, where $n=\lfloor s_2/(s_1-s_2)\rfloor$. We prove both
assertions. We exhibit, in closed form, the pure subgame perfect equilibrium
outcome at \emph{every} heap below that value---the same in both
conventions---and evaluate both conventions at the value itself, obtaining
$o_{\FvF}=(s_2+ns_1,\,(n+1)s_2)$ and $o_{\AvA}=(s_2+ns_1,\,ns_1)$. Existence
and minimality are then immediate, the mover's coordinate agrees at the
distinguishing heap, and the whole discrepancy is
$-\bigl(s_2\bmod(s_1-s_2)\bigr)$ in the non-mover's coordinate. We claim no
priority: a manuscript of Kulkarni and Larsson listed as in preparation in that
paper is on exactly this subject and could not be obtained.
\end{abstract}

\maketitle

\section{The game and the conjecture}

Fix a finite $S\subset\mathbb Z_{>0}$. In the \emph{self-interest cumulative
subtraction game} on $S$ two players alternate; from a heap of $h$ tokens the
mover removes some $s\in S$ with $s\le h$ and adds $s$ to their \emph{own}
tally; play ends when $h<\min S$. The game is not zero-sum, and the relevant
solution concept is a pure subgame perfect equilibrium (PSPE). Write
$o(h)=(o^1(h),o^2(h))$ for the PSPE payoff pair, $o^1$ being the mover's total.
Then $o(h)=(0,0)$ for $h<\min S$ and
\begin{equation}\label{eq:rec}
  o^1(h)=\max_{s\in S,\ s\le h}\bigl(s+o^2(h-s)\bigr),
  \qquad
  o^2(h)=o^1(h-s^*),
\end{equation}
where $s^*$ is the maximiser. When two moves are \emph{indifferent}, i.e.\ they
attain the same value of $s+o^2(h-s)$, \eqref{eq:rec} does not determine
$o^2(h)$ and a convention is needed. Following
Bhagat--Kulkarni--Larsson--Murali \cite{BKLM}, in the \emph{friendly}
convention $\FvF$ the mover selects, among indifferent moves, one
\emph{maximising} $o^1(h-s)$, and in the \emph{antagonistic} convention $\AvA$
one \emph{minimising} it. Both recursions are well defined and single-valued;
we write $o_{\FvF}$ and $o_{\AvA}$, and set
\[
  \delta^1(h)=o^1_{\AvA}(h)-o^1_{\FvF}(h),
  \qquad
  \delta^2(h)=o^2_{\AvA}(h)-o^2_{\FvF}(h).
\]
Call $h$ a \emph{discrepancy heap} if $o_{\FvF}(h)\neq o_{\AvA}(h)$.

From here on $S=\{s_2,s_1\}$ with $0<s_2<s_1$, and it is convenient to put
\begin{equation}\label{eq:abd}
  a:=s_2,\qquad b:=s_1,\qquad d:=b-a>0 .
\end{equation}

The object of this note is the following conjecture of \cite{BKLM}, which we
quote verbatim, translating only its internal labels into words.

\begin{conjecture}[\cite{BKLM}]\label{con:halfline}
Consider $S=\{s_2,s_1\}$ in the non-dominant regime where $s_1<2s_2$, and
suppose, for all $k\in\mathbb N$, $s_1/s_2\neq(k+1)/k$. Then there exists a
heap $h$ that distinguishes the outcomes of the $\FvF$ and $\AvA$ conventions.
The smallest such heap is defined by
\begin{equation}\label{eq:smallestheap}
  h=(n+2)s_2+ns_1 .
\end{equation}
\end{conjecture}

The index $n$ is fixed in \cite{BKLM} by the interval condition
$\frac{n+1}{n}>\frac{s_1}{s_2}>\frac{n+2}{n+1}$, and a remark there gives the
equivalent closed form $n=\lfloor s_2/(s_1-s_2)\rfloor$. The two agree, and
their agreement already encodes the hypotheses:

\begin{lemma}\label{lem:hyp}
With the notation \eqref{eq:abd}, the hypotheses of
Conjecture~\ref{con:halfline} are exactly $0<d<a$ together with $d\nmid a$; and
under them $n=\lfloor a/d\rfloor\ge1$ is the unique integer with
$n<a/d<n+1$.
\end{lemma}

\begin{proof}
$s_1/s_2=1+d/a$, so $\frac{n+1}{n}>1+\frac da$ reads $a/d>n$ and
$1+\frac da>\frac{n+2}{n+1}$ reads $a/d<n+1$. Non-dominance $s_1<2s_2$ is
$d<a$, i.e.\ $a/d>1$, i.e.\ $n\ge1$. The excluded ratios are
$s_1/s_2=(k+1)/k$, i.e.\ $1+d/a=1+1/k$, i.e.\ $a=kd$; so excluding them is
exactly $d\nmid a$, which is what makes both inequalities of the interval
condition strict.
\end{proof}

Throughout, therefore,
\begin{equation}\label{eq:data}
  0<d<a,\quad d\nmid a,\quad
  n:=\Bigl\lfloor\frac ad\Bigr\rfloor\ge1,\quad
  \varepsilon:=a-nd\in(0,d),\quad
  \alpha_t:=a-td,
\end{equation}
and the heap named by \eqref{eq:smallestheap} is
\begin{equation}\label{eq:H}
  H:=(n+2)a+nb=(2n+2)a+nd .
\end{equation}

The neighbouring regime is already settled in \cite{BKLM}: a proposition proved
there forbids any discrepancy once $s_1\ge2s_2$. For
Conjecture~\ref{con:halfline} \cite{BKLM} offer computational evidence, a figure
reaching $n=48$ and further simulations.

\section{The result}

\begin{theorem}\label{thm:main}
Assume \eqref{eq:data} and let $H$ be as in \eqref{eq:H}. Call a pair of
integers $(m,i)$ \emph{admissible} if $m\ge0$ and $0\le i\le\lceil m/2\rceil$,
and put
\[
  v(m,i):=ia+(m-i)b=mb-id .
\]
Let $W:=\{v(m,i):(m,i)\ \text{admissible}\}$, let
$V(h):=\max\{w\in W:w\le h\}$, let $(m,i)$ be the admissible coordinates of
$V(h)$ \textup{(}unique for $h\le H$, by Lemma~\textup{\ref{lem:lattice}}\/\textup{)}
and $t:=\lfloor m/2\rfloor$. Then for every $h<H$, in \emph{both}
conventions,
\begin{equation}\label{eq:inv}
  \bigl(o^1(h),o^2(h)\bigr)=
  \begin{cases}
    \bigl(tb,\ tb-id\bigr), & m=2t,\\[2pt]
    \bigl((t+1)b-id,\ tb\bigr), & m=2t+1,
  \end{cases}
\end{equation}
while at $h=H$
\begin{equation}\label{eq:atH}
  o_{\FvF}(H)=\bigl(a+nb,\ (n+1)a\bigr)
  \neq
  \bigl(a+nb,\ nb\bigr)=o_{\AvA}(H).
\end{equation}
\end{theorem}

\begin{corollary}\label{cor:conj}
Conjecture~\ref{con:halfline} is true. A discrepancy heap exists, the smallest
one is exactly $H=(n+2)s_2+ns_1$, and there
$\delta^1(H)=0$ while
\[
  \delta^2(H)=-\varepsilon=-\bigl(s_2\bmod(s_1-s_2)\bigr)<0 .
\]
\end{corollary}

\begin{proof}
By \eqref{eq:inv} the outcome at every $h<H$ is the same in both conventions,
so no $h<H$ is a discrepancy heap; by \eqref{eq:atH}, $H$ is one. The two
first coordinates in \eqref{eq:atH} are equal, and
$nb-(n+1)a=n(a+d)-(n+1)a=nd-a=-\varepsilon$.
\end{proof}

Equation \eqref{eq:inv} is a complete description of the equilibrium payoff and
aggregate move composition below $H$: PSPE play lasts $m$ moves and consumes $i$
copies of $s_2$ and $m-i$ copies of $s_1$, leaving the residue $h-V(h)$ on the
heap. It does not single out an equilibrium strategy or a unique move sequence.

\section{The lattice $W$}

Write $L_m:=\{v(m,i):0\le i\le\lceil m/2\rceil\}$, an arithmetic progression of
step $d$ with
\[
  \max L_m=mb,\qquad
  \min L_{2t}=t(a+b),\qquad
  \min L_{2t+1}=(t+1)a+tb .
\]

\begin{lemma}\label{lem:lattice}
Assume \eqref{eq:data}. Then
\begin{enumerate}
\item[\textup{(F1)}]
  $\min L_{2t+1}-\max L_{2t}=\alpha_t$ and
  $\min L_{2t+2}-\max L_{2t+1}=\alpha_t$ for every $t\ge0$;
\item[\textup{(F2)}]
  $\alpha_t\ge\alpha_n=\varepsilon>0$ for $t\le n$, hence
  $L_0,L_1,\dots,L_{2n+2}$ are pairwise disjoint and increasing, and every
  $w\in W$ with $w\le H$ has a \emph{unique} admissible coordinate pair;
\item[\textup{(F3)}]
  $\min L_{2n+2}=H+d$ and $\max L_{2n+1}=H+(d-\varepsilon)>H$; consequently
  $W\cap[0,H]\subseteq L_0\cup\dots\cup L_{2n+1}$, $H\notin W$, and
  $\max\bigl(W\cap[0,H]\bigr)=v(2n+1,1)=2nb+a=H-\varepsilon$;
\item[\textup{(F4)}]
  for admissible $(m,i)$ with $m\le2n+1$ --- which by \textup{(F3)} covers every
  point of $W\cap[0,H]$ --- the gap in $W$ immediately above $v(m,i)$ is $d$ if
  $i\ge1$ and $\alpha_{\lfloor m/2\rfloor}$ if $i=0$.
\end{enumerate}
\end{lemma}

\begin{proof}
(F1) $\min L_{2t+1}-\max L_{2t}=(t+1)a+tb-2tb=(t+1)a-tb=a-td=\alpha_t$, using
$b=a+d$; and
$\min L_{2t+2}-\max L_{2t+1}=(t+1)(a+b)-(2t+1)b=(t+1)a-tb=\alpha_t$.

(F2) $\alpha_t=a-td$ is decreasing in $t$ and $\alpha_n=a-nd=\varepsilon>0$ by
\eqref{eq:data}; so all the gaps in (F1) with $t\le n$ are strictly positive,
each $L_m$ has step $d>0$, and $L_0<L_1<\dots<L_{2n+2}$ as sets of integers.
Uniqueness of coordinates on $[0,H]$ follows, because by (F3) no
$v(m,i)$ with $m\ge2n+2$ lies in $[0,H]$. (Positivity of $\varepsilon$ is
exactly where $d\nmid a$ enters.)

(F3) $\min L_{2n+2}=(n+1)(a+b)=(2n+2)a+(n+1)d=H+d$, and
$\max L_{2n+1}=(2n+1)b=(2n+1)a+(2n+1)d=H+\bigl((n+1)d-a\bigr)=H+(d-\varepsilon)$,
which exceeds $H$ since $\varepsilon<d$. As $\min L_m$ is increasing in $m$,
every $v(m,i)$ with $m\ge2n+2$ is at least $H+d$, which gives the inclusion.
Inside $L_{2n+1}$, $v(2n+1,i)\le H$ reads $(1-i)d\le\varepsilon$, true for
$i\ge1$ and false for $i=0$; so the largest element of $W$ below $H$ is
$v(2n+1,1)=(2n+1)b-d=2nb+a$, at distance
$H-(2nb+a)=(2n+2)a+nd-(2n+1)a-2nd=\varepsilon$ from $H$. Since
$0<\varepsilon<d$, that distance is not a multiple of $d$, so $H\notin L_{2n+1}$;
and $H<\min L_{2n+2}$ while $H>\max L_{2n}=2nb$, so $H\notin W$.

(F4) Let $m\le2n+1$. Since $\min L_{m'}$ is increasing in $m'$, every element of
$W$ lying outside $L_0\cup\dots\cup L_{2n+2}$ is at least
$\min L_{2n+3}>\min L_{2n+2}$; so on $\bigl[0,\min L_{2n+2}\bigr]$ the elements
of $W$ are exactly those of $L_0\cup\dots\cup L_{2n+2}$, ordered as in (F2).
Both $v(m,i)$ and $\min L_{m+1}$ lie in that range, since
$\max L_{2n+1}<\min L_{2n+2}$ by (F3). If $i\ge1$ then the next element of
$L_m$, and hence of $W$, is $v(m,i-1)=v(m,i)+d$. If $i=0$ then
$v(m,0)=\max L_m$ and the next element of $W$ is $\min L_{m+1}$, at distance
$\alpha_{\lfloor m/2\rfloor}$ by (F1).
\end{proof}

\section{Proof of Theorem~\ref{thm:main}}

We induct on $h$. For $h<a$ the only admissible value $\le h$ is $v(0,0)=0$, so
$m=i=t=0$ and \eqref{eq:inv} reads $o(h)=(0,0)$, which is the base of
\eqref{eq:rec}.

Let $a\le h\le H$ and let $(m,i)$ be the coordinates of $V(h)$, $t=\lfloor
m/2\rfloor$, $r:=h-V(h)$. Since $v(1,1)=a\le h$ we have $m\ge1$, and by (F3)
$m\le2n+1$, hence
\begin{equation}\label{eq:tn}
  t\le n,\qquad\text{so}\qquad \alpha_t\ge\varepsilon>0 .
\end{equation}
Put $c_a:=a+o^2(h-a)$ and, when $h\ge b$, $c_b:=b+o^2(h-b)$; by \eqref{eq:rec}
the mover takes the larger, and on a tie $\FvF$ takes the $s$ maximising
$o^1(h-s)$ while $\AvA$ takes the $s$ minimising it. By admissibility
$0\le i\le t$ when $m=2t$ and $0\le i\le t+1$ when $m=2t+1$, so exactly one of
the following five cases holds; the split is therefore exhaustive and disjoint.
In each case we use the inductive hypothesis \eqref{eq:inv} at $h-a$ and $h-b$,
whose coordinates we first locate using (F4).

\medskip
\noindent\textbf{Case 1: $m=2t$, $1\le i\le t$.}
Here $r<d$ by (F4). Since $b-a=d$ and $v(m,i)-b=v(2t-1,i)$,
$v(m,i)-a=v(2t-1,i-1)$, and since the gaps above $v(2t-1,i)$ and
$v(2t-1,i-1)$ are $d$ (or $\alpha_{t-1}=\alpha_t+d>d$ when $i=1$) we get
$V(h-b)=v(2t-1,i)$ and $V(h-a)=v(2t-1,i-1)$. Both have $m'=2t-1$ odd with
$t'=t-1$, so \eqref{eq:inv} gives $o^2(h-b)=o^2(h-a)=(t-1)b$ and
\[
  c_b-c_a=b-a=d>0 .
\]
Thus $s^*=s_1$ \emph{uniquely}, $o^1(h)=c_b=tb$ and
$o^2(h)=o^1(h-b)=tb-id$: this is \eqref{eq:inv} for $m=2t$, in both
conventions.

\medskip
\noindent\textbf{Case 2: $m=2t$, $i=0$, $t\ge1$.}
Here $r<\alpha_t$ by (F4). Now $h-b=v(2t-1,0)+r$ and
$h-a=v(2t-1,0)+d+r$, and $d+r<d+\alpha_t=\alpha_{t-1}$, so \emph{both} residues
sit on $v(2t-1,0)$ and $o^2(h-a)=o^2(h-b)=(t-1)b$. Again $c_b-c_a=d>0$, so
$s^*=s_1$ uniquely, $o^1(h)=tb$ and $o^2(h)=o^1(h-b)=tb$, which is
\eqref{eq:inv} with $i=0$.

\medskip
\noindent\textbf{Case 3: $m=2t+1$, $1\le i\le t$ \textup{(}so $t\ge1$\textup{)}.}
Here $r<d$. As $v(m,i)-b=v(2t,i)$ with $i\ge1$, we get $V(h-b)=v(2t,i)$,
$o^2(h-b)=tb-id$, $o^1(h-b)=tb$ and
\[
  c_b=(t+1)b-id .
\]
Also $v(m,i)-a=v(2t,i-1)$ exactly. Two sub-cases, according to (F4) applied at
$v(2t,i-1)$, whose gap above is $d$ if $i\ge2$ and $\alpha_t$ if $i=1$.

\emph{Case 3A: $i\ge2$, or $i=1$ and $r<\alpha_t$.} Then
$V(h-a)=v(2t,i-1)$, so $o^2(h-a)=tb-(i-1)d$ and
$c_a=a+tb-(i-1)d=(t+1)b-id=c_b$: the two moves are indifferent. But the
tie-breaking comparands agree, $o^1(h-a)=tb=o^1(h-b)$ by \eqref{eq:inv}, so
both conventions return the same pair
$\bigl((t+1)b-id,\ tb\bigr)$ --- the tie is not load-bearing.

\emph{Case 3B: $i=1$ and $r\ge\alpha_t$.} Since $r<d$ this forces
$\alpha_t<d$, i.e.\ $a<(t+1)d$, i.e.\ $t\ge n$; with \eqref{eq:tn}, $t=n$.
Hence $\alpha_t=\varepsilon$, $V(h)=v(2n+1,1)=H-\varepsilon$ by (F3), and
$r\ge\varepsilon$ forces $h\ge H$, so $h=H$. \emph{Case 3B therefore occurs at
$h=H$ and nowhere below it}; this is the minimality assertion. At $h=H$ we have
$r=\varepsilon$ and
\[
  H-a=(2n+1)a+nd=v(2n+1,n+1),\qquad
  H-b=v(2n,1)+\varepsilon ,
\]
so by \eqref{eq:inv}
\[
  o^1(H-a)=(n+1)b-(n+1)d=(n+1)a,\quad o^2(H-a)=nb,
\]
\[
  o^1(H-b)=nb,\quad o^2(H-b)=nb-d .
\]
Hence $c_a=a+nb$ and $c_b=b+nb-d=(n+1)a+nd=a+nb=c_a$: again a tie. This time
the comparands \emph{split},
\[
  o^1(H-a)-o^1(H-b)=(n+1)a-nb=a-nd=\varepsilon>0 .
\]
So $\FvF$, which maximises the opponent's payoff, plays $s_2$ and the non-mover
receives $o^1(H-a)=(n+1)a$; $\AvA$ plays $s_1$ and the non-mover receives
$o^1(H-b)=nb$. Both give $o^1(H)=a+nb$. This is \eqref{eq:atH}.

\medskip
\noindent\textbf{Case 4: $m=2t+1$, $i=0$.}
Here $r<\alpha_t$ by (F4), and $t\le n-1$: indeed $v(2n+1,0)>H$ by (F3). Now
$h-b=v(2t,0)+r$ with $r<\alpha_t$, so $V(h-b)=v(2t,0)$, $o^2(h-b)=tb$ and
$c_b=(t+1)b$. For $h-a=v(2t,0)+d+r$ there are two possibilities --- if
$d+r<\alpha_t$ then $V(h-a)=v(2t,0)$, and otherwise
$v(2t,0)+\alpha_t\le h-a<v(2t,0)+\alpha_t+d$ so $V(h-a)=v(2t+1,t+1)$ --- but
\eqref{eq:inv} returns $o^2(h-a)=tb$ in \emph{either} case, so $c_a=a+tb$ and
\[
  c_b-c_a=b-a=d>0 .
\]
Thus $s^*=s_1$ uniquely, $o^1(h)=(t+1)b$ and $o^2(h)=o^1(h-b)=tb$, which is
\eqref{eq:inv}.

\medskip
\noindent\textbf{Case 5: $m=2t+1$, $i=t+1$.}
Then $V(h)=\min L_{2t+1}=(t+1)a+tb$ and $r<d$. If $t=0$ then $a\le h<b$ and
only $s_2$ is legal, giving $o(h)=(a,0)$, which is \eqref{eq:inv}. If $t\ge1$,
then $h-a=v(2t,t)+r$ with $t\ge1$, so $V(h-a)=v(2t,t)$ and
$o^2(h-a)=tb-td=ta$, whence $c_a=(t+1)a$; and
$h-b=v(2t-1,0)+\alpha_t+r$ with $\alpha_t+r<\alpha_t+d=\alpha_{t-1}$, so
$V(h-b)=v(2t-1,0)$ and $o^2(h-b)=(t-1)b$, whence $c_b=tb$. Therefore
\[
  c_a-c_b=(t+1)a-tb=\alpha_t>0
\]
by \eqref{eq:tn}: the \emph{sacrifice} $s^*=s_2$ is uniquely optimal. Then
$o^1(h)=(t+1)a=(t+1)b-(t+1)d$ and $o^2(h)=o^1(h-a)=tb$, which is
\eqref{eq:inv} with $i=t+1$.

\medskip
Every case with $h<H$ --- that is, Cases 1, 2, 3A, 4, 5 --- produces a
convention-independent pair satisfying \eqref{eq:inv}, and $h=H$ falls in
Case 3B and produces \eqref{eq:atH}. This completes the induction and the
proof. \qed

\begin{remark}[where the hypotheses are used]
The two hypotheses enter together, through $\varepsilon>0$: non-dominance
$d<a$ gives $n\ge1$, and $d\nmid a$ then gives $\varepsilon=a-nd>0$, hence
$\alpha_t\ge\varepsilon>0$ for $t\le n$ in \eqref{eq:tn}, hence the strict
sacrifice of Case 5 and the strict comparand gap $\varepsilon$ at $H$. At
$d\ge a$ the ordering of Lemma~\ref{lem:lattice} collapses and the proposition
of \cite{BKLM} recalled in \S1 takes over; at $\varepsilon=0$ Case 3B
degenerates into Case 3A and that tie is not load-bearing.
\end{remark}

\section{An example}

The hypothesis set $\{(a,d):1\le d<a,\ d\nmid a\}$ is non-empty from $a=3$,
$d=2$, i.e.\ $S=\{3,5\}$, which is the running example of \cite{BKLM}. Here
$n=1$, $\varepsilon=1$ and $H=3\cdot3+1\cdot5=14$. The admissible values
$\le13$ are $0,3,5,8,10,11,13$ with coordinates
$(0,0),(1,1),(1,0),(2,1),(2,0),(3,2),(3,1)$, and \eqref{eq:inv} gives the
block table

\begin{center}
\begin{tabular}{l|ccccccc}
$h$ & $[0,3)$ & $[3,5)$ & $[5,8)$ & $[8,10)$ & $[10,11)$ & $[11,13)$ & $[13,14)$\\
\hline
$(o^1,o^2)$ & $(0,0)$ & $(3,0)$ & $(5,0)$ & $(5,3)$ & $(5,5)$ & $(6,5)$ & $(8,5)$
\end{tabular}
\end{center}

\noindent
while \eqref{eq:atH} gives $o_{\FvF}(14)=(8,6)$ and $o_{\AvA}(14)=(8,5)$. The
table for $S=\{3,5\}$ in \cite{BKLM} prints the columns $h=0,\dots,15$ with
$\delta^1\equiv0$ and the single non-zero entry $\delta^2(14)=-1$, and records
the PSPE move sequences $3\text{-}3\text{-}5\text{-}3$ under $\FvF$ and
$5\text{-}5\text{-}3$ under $\AvA$, which is Case 3B's $s_2$ against $s_1$.

\section{Scope and verification}

What is proved is Conjecture~\ref{con:halfline} in full, uniformly in $n\ge1$.
Not settled: the companion statement of \cite{BKLM}, that at the excluded ratios
$s_1/s_2=(k+1)/k$ there is no discrepancy at any heap---nothing here bears on
it, and at $d\mid a$ one has $\alpha_n=0$, Lemma~\ref{lem:lattice}(F2) fails at
$t=n$ and the coordinates of Theorem~\ref{thm:main} are not even well defined,
so in particular this proof does not show the absence of a discrepancy at or
below $H$ there; heaps $h>H$, where the outcome genuinely is
convention-dependent and no claim is made; subtraction sets with $|S|\ge3$; and
formalisation, this being a hand proof, not checked in a proof assistant.
Reference~[12] of \cite{BKLM} is \cite{KL}, by two of its four authors and on
this subject, listed there as \emph{in preparation}, with no arXiv identifier
and no DOI; we were unable to obtain it, so we do not know what it contains. We
claim no novelty or priority against it, nor against \cite{CLMW}, whose
solution of the two-action \emph{zero-sum} game already contains part of the
structure of play below $H$ (a zero-sum game has a single-valued outcome and so
no conventions to disagree, so \cite{CLMW} states neither half of
Conjecture~\ref{con:halfline}), nor against \cite{LMZ}, where the instance
$S=\{3,5\}$ at $h=14$ is on record.

A program accompanying this note, \texttt{verify.py}, with a transcript of one
run in \texttt{verify.output.txt}, brute-forces \eqref{eq:rec} in both
conventions in exact integer arithmetic over the $4568$ pairs with $s_2\le100$
satisfying \eqref{eq:data}, and checks the invariant \eqref{eq:inv} and its
convention-independence at every $h<H$, the values \eqref{eq:atH} and
$\delta^2(H)=-\varepsilon$, that the least discrepancy heap is exactly $H$,
each part of Lemma~\ref{lem:lattice}, and the exhaustiveness, disjointness and
per-case arithmetic of the five-case split of \S4; all $32$ of its checks pass.
Since $d\nmid a$ forces $d\ge2$, that box reaches only $n\le49$, so for the
formula itself it adds little to the computational evidence already in
\cite{BKLM}; what it adds is the check of the structural claims. It audits the
induction on a finite range and does not certify Theorem~\ref{thm:main}, which
is proved above by hand.

\begin{thebibliography}{9}

\bibitem{BKLM}
A.~Bhagat, T.~Kulkarni, U.~Larsson, and A.~Murali,
\emph{Tie-breaking in self interest cumulative subtraction games},
arXiv:2510.24280, 2025 (v1, 28 Oct 2025; v2, 20 Jan 2026).
\href{https://arxiv.org/abs/2510.24280}{arXiv:2510.24280}.
Quotations above are from the version~2 e-print.

\bibitem{CLMW}
G.~Cohensius, U.~Larsson, R.~Meir, and D.~Wahlstedt,
\emph{Cumulative subtraction games},
Electron. J. Combin. \textbf{26} (2019), no.~4, Paper 4.52.
\href{https://doi.org/10.37236/7904}{doi:10.37236/7904};
\href{https://arxiv.org/abs/1805.09368}{arXiv:1805.09368}.

\bibitem{LMZ}
U.~Larsson, R.~Meir, and Y.~Zick,
\emph{Cumulative games: who is the current player?},
arXiv:2005.06326, 2020.
\href{https://arxiv.org/abs/2005.06326}{arXiv:2005.06326}.
The observation cited above is in the December 2025 revision.

\bibitem{KL}
T.~Kulkarni and U.~Larsson,
\emph{Discrepancies in various 2-player self-interest cumulative subtraction
games}, in preparation. Cited as reference~[12] of \cite{BKLM}.

\end{thebibliography}

\end{document}
